\documentclass[11pt]{article}
\usepackage[margin=1in]{geometry}
\usepackage{amsmath,amssymb,amsfonts,mathtools,bm}
\usepackage[T1]{fontenc}
\usepackage[utf8]{inputenc}
\title{Aligned Mathematical Structures}
\author{Source-backed grouped formula sample}
\date{}
\begin{document}
\maketitle
\section{Expressions}
\paragraph{Expression 1.} The following expression is taken from a source-backed formula corpus.

\begin{align*}\begin{array}{ccl}\xi _k^\alpha (x) & = & \left( \left( \frac 1{R\left( A(x)\right) R\left(A(x)\right) }\right) ^{1/4}\right) _{kk^{^{\prime }}}^{\alpha \alpha^{^{\prime }}}B_{k^{^{\prime }}}^{\alpha ^{^{\prime }}}(x) \\ & = & \frac 1{\pi \sqrt{2}}\int_0^\infty d\lambda \lambda ^{-1/4}\frac1{\lambda +R\bullet R}\bullet B(x) \end{array}\end{align*}

\paragraph{Expression 2.} The following expression is taken from a source-backed formula corpus.

\begin{align*}s&=\sum_{i=0}^{r-1}\left\{\langle\frac{-2np^i}{q-1}\rangle+\langle\frac{ndp^i}{q-1}\rangle+\langle\frac{-n(d-1)p^i}{q-1}\rangle-\langle\frac{-np^i}{q-1}\rangle\right\}\\&=\sum_{i=0}^{r-1}\left\{\lfloor\frac{-np^i}{q-1}\rfloor-\lfloor\frac{-2np^i}{q-1}\rfloor-\lfloor\frac{ndp^i}{q-1}\rfloor-\lfloor\frac{-n(d-1)p^i}{q-1}\rfloor\right\},\end{align*}

\paragraph{Expression 3.} The following expression is taken from a source-backed formula corpus.

\begin{align*}\frac{\partial\eta_k(\xi)}{\partial \xi} =&\frac{\sigma_k}{\xi^2}+\frac{2\log(C_k)}{\xi^3}+\frac{k}{\xi^3}+\frac{k}{\xi^2\sigma_k}\\&-\left(\frac{1}{2\sigma_k}+\frac{2k\log(\sigma_k+\xi)}{\xi^3}+\frac{\log\left(\frac{1}{2}\left(1+\frac{\xi}{\sigma_k}\right)\right)}{\xi^3}+\frac{\xi^2-\sigma_k^2}{2\xi^2\sigma_k^2(\sigma_k+\xi)}\right).\end{align*}

\paragraph{Expression 4.} The following expression is taken from a source-backed formula corpus.

\begin{align*}S_{\omega,\mathbf{c}}(U)&:=E(U)+\omega Q(U)+\mathbf{c}\cdot \mathbf{P}(U),\\K_{\omega,\mathbf{c}}(U)&:=\partial_{\lambda}S_{\omega,\mathbf{c}}(\lambda U)|_{\lambda =1}=2L(U)+3N(U)+2\omega Q(U)+2 \mathbf{c}\cdot \mathbf{P}(U),\\L_{\omega,\mathbf{c}}(U)&:=K_{\omega, \mathbf{c}}(U)-3N(U)=2L(U)+2\omega Q(U)+2\mathbf{c}\cdot \mathbf{P}(U).\end{align*}

\paragraph{Expression 5.} The following expression is taken from a source-backed formula corpus.

\begin{align*}U^{\prime}_{[4]\times[N]} = - \left( \! \begin{array}{cccc} U^{\prime}_{11} & U^{\prime}_{12} & \cdots & U^{\prime}_{1N} \\U^{\prime}_{21} & U^{\prime}_{22} & \cdots & U^{\prime}_{2N} \\U^{\prime}_{31} & U^{\prime}_{32} & \cdots & U^{\prime}_{3N} \\U^{\prime}_{41} & U^{\prime}_{42} & \cdots & U^{\prime}_{4N} \\\end{array} \! \right),\end{align*}
\end{document}
